
In standard MPC, stability and recursive feasibility are guaranteed
by terminal sets $\mathcal{X}_f$ and terminal costs $l_f$.
IL-MPC constructs these \textit{ingredients} dynamically
from previous trajectories:

\textbf{Sampled Safe Set} ($\mathcal{SS}^e$)
The collection of all successful
state trajectories recorded up to episode $e$.

\textbf{Convex Safe Set} ($\mathcal{CS}^e$)
The convex hull of $SS^e$.
For linear systems,
any convex combination of successful trajectories
is also a feasible trajectory.
This set is \textbf{control invariant}.

\textbf{Learned Terminal Cost} ($P^e(x)$)
Defined as a barycentric function over the safe set.
It assigns each point the minimum \textit{cost-to-go}
based on a convex combination of previous costs.

\textbf{Guarantees}
Under linear dynamics, the resulting MPC is recursively feasible,
asymptotically stable, and the closed-loop cost
is non-increasing across episodes, eventually converging to the
infinite-horizon optimal solution.

\subsection{Global Tuning Approaches (Bayesian Methods)}

Used when the relationship between cost parameters (e.g., $Q, R, P$)
and a high-level performance metric (e.g., lap time) is unknown or
non-analytical.

\textbf{Bayesian Optimization (BO)}
Uses a \textbf{Gaussian Process} as a surrogate model
to represent the unknown performance function.
An \textbf{Acquisition Function} (like Lower Confidence Bound - LCB)
suggests the next parameters to test by balancing
exploration (uncertainty) and exploitation (predicted performance).

\textbf{Safe BO}
Restricts exploration to parameters that are likely
to maintain performance above a minimum safety threshold.

\textbf{Contextual BO}
Adapts the tuning to changing environments or system parameters
(e.g., racing different cars)
by treating model mismatches as a \textit{context} variable $\theta$.

\textbf{Bayesian MPC}
Uses \textbf{Thompson Sampling} to sample parameters
from a posterior distribution and use them in the MPC,
minimizing \textit{cumulative regret}.

\subsection{Local Tuning Approaches (Differentiable MPC)}

If the performance metric is differentiable,
we can use gradient-based methods.

\textbf{Differentiable MPC}
Treats the entire MPC solver as a differentiable layer.
This allows the computation of gradients of the
optimal trajectories $\frac{\partial u^*}{\partial \theta}$
with respect to cost or dynamics parameters.

\textbf{RL-MPC}
Uses Reinforcement Learning to update cost parameters $\theta$
to minimize long-term cost.

\textbf{ZipMPC}
A student-teacher framework.
A \textit{Student} (short-horizon MPC) is trained via imitation learning
to mimic the behavior of a \textit{Teacher}
(computationally expensive long-horizon MPC)

\subsection{Approximate Explicit MPC}

Standard MPC requires solving an optimization problem online, which
may be too slow for high-frequency applications. Explicit MPC moves
the computation offline.

\subsubsection{Explicit MPC for Linear Systems}

For linear systems with polyhedral constraints,
the MPC solution is a \textbf{Multiparametric Quadratic Program (mp-QP)}.

The state space is partitioned into polyhedral regions.
In each region,
the control law is a unique \textbf{Piecewise Affine (PWA)} function:
$\pi^*(x) = K_i x + k_i$.

\textbf{Pros/Cons}
Online evaluation is just a \textit{point-in-polyhedron} search (very fast),
but the number of regions grows exponentially with
constraints and horizon (high memory usage).

\subsubsection{Approximate Explicit MPC for Nonlinear Systems}

Nonlinear systems generally lack an analytical explicit solution.
Instead, we learn a function approximation $\tilde{\pi}(x) \approx \pi^*(x)$.

\textbf{Regression Techniques}
Neural Networks (NN) or Kernel Interpolation are used
to fit a model to a dataset of offline-solved MPC problems.

\textbf{Verification and Safety}
Because approximation errors can break stability/safety,
two approaches are used:

\begin{enumerate}
  \item \textbf{Safety Filters}
    Project the approximate input $\tilde{u}$ onto a control invariant set
    to ensure constraints are always met.
  \item \textbf{Probabilistic Verification}
    Use the \textbf{Hoeffding Bound} to certify
    the controller’s success rate $\mu$ with a specific
    confidence level $(1-\delta_h)$ based on a finite number of test samples.
\end{enumerate}

